\printedpage{9}
% head: A. Grothendieck : Dérivateurs, Ch. XVII, Édition des Universités de Montpellier II et Paris VII
% foot: Version 9 février 2019
\[
\begin{tikzcd}
X_J \arrow[r] \arrow[d, "u_J"'] \arrow[dr, phantom, "\cocart"] & X_K \arrow[r] \arrow[d, "\tilde{u}_{J,K}"'] \arrow[dr, phantom, "\cocart"] & X_L \arrow[d, "\tilde{u}_{J,L}"'] \arrow[dd, bend left=40, "\tilde{u}_{K,L}"] \\
Y_J \arrow[r] & Y_{J,K} \arrow[r] \arrow[d, "u_{J,K}"'] \arrow[dr, phantom, "\cocart"] & Y_{J,L} \arrow[d, "u^L_{J,K}"'] \arrow[dd, bend left=40, "u_{J,L}"] \\
 & Y_K \arrow[r] & Y_{K,L} \arrow[d, "u_{K,L}"'] \\
 & & Y_L\ ,
\end{tikzcd}
\tag*{(2)}
\]
giving rise to a new arrow $u^L_{J,K}$, deduced from $u_{J,K}$ by the cobase change \msp{page 15} $Y_{J,K} \longrightarrow Y_{J,L}$. We shall use this diagram in the two cases where $(J \subset K \subset L)$ is one of the two triples
\[
\left\{
\begin{aligned}
& I'_0 \hookrightarrow I_0 \hookrightarrow I \\
& I'_0 \hookrightarrow I' \hookrightarrow I,
\end{aligned}
\right.
\]
arising from the cocartesian square (in Cat) (MV square)
\[
\begin{tikzcd}
I'_0 \arrow[r, hook] \arrow[d, hook] & I_0 \arrow[d, hook] \\
I' \arrow[r, hook] & I\ .
\end{tikzcd}
\tag*{(3)}
\]

We must insert the two corresponding diagrams of type $(*)$ \texttt{[on page 10]} into one and the same diagram. We shall note in this connection that the squares in $\mathcal{M}$ deduced from (3) by applying to it $J \longmapsto X_J$ and $J \longmapsto Y_J$
\[
\begin{matrix} (4_X) \\[12ex] (4_Y) \end{matrix}
\qquad\qquad
\begin{tikzcd}[column sep=small, row sep=small]
X_{I'_0} \arrow[rr] \arrow[dr] \arrow[dd, dotted, "u_{I'_0}"] & & X_{I_0} \arrow[dr] \arrow[dd, dotted, "u_{I_0}"] & \\
 & X_{I'} \arrow[rr] \arrow[dd, dotted, "u_{I'}"'] & & X_I \arrow[dd, dotted, "u_I"'] \\
Y_{I'_0} \arrow[rr] \arrow[dr] & & Y_{I_0} \arrow[dr] & \\
 & Y_{I'} \arrow[rr] & & Y_I
\end{tikzcd}
\qquad\qquad
\left.\vrule height 9ex depth 9ex width 0pt\right\}\ (4)
\]
are cocartesian in $\mathcal{M}$. (Transitivity of the $\varinjlim$ \ldots). % [sic] cocartésien -> cocartésiens, agreement with « les carrés »

\msp{page 16}
We have a homomorphism $(u_{I'_0}, u_{I_0}, u_{I'}, u_I)$ deduced from $u$ from the square $(4_X)$ into the square $(4_Y)$, % [sic] les carré -> le carré, singular noun with plural article
and this homomorphism too fits into a diagram, including the two avatars of (2) already mentioned. Here is the relevant diagram: % [sic] signalé -> signalés, agreement with « les deux avatars »
\printedpage{10}
% head: A. Grothendieck : Dérivateurs, Ch. XVII, Édition des Universités de Montpellier II et Paris VII
% foot: Version 9 février 2019
% the footnote mark (9) stands in the printed label of the display; \footnotemark/\footnotetext used because \footnote cannot be set inside a display tag
\[
\begin{tikzcd}[column sep=small, row sep=small]
 & X_{I_0} \arrow[rrr] \arrow[dd] & & & X_I \arrow[dd] \\
X_{I'_0} \arrow[ur] \arrow[rr] \arrow[dd] & & X_{I'} \arrow[urr] \arrow[dd] & & \\
 & Y_{I'_0,I_0} \arrow[rrr] \arrow[dd, "u_{I'_0,I_0}"] & & & Y_{I'_0,I} \arrow[dd, "u^I_{I'_0,I_0}"] \arrow[dddl] \\
Y_{I'_0} \arrow[ur] \arrow[rr] \arrow[dr, dotted] \arrow[dddrr, dotted] & & Y_{I'_0,I'} \arrow[urr] \arrow[ddd] & & \\
 & Y_{I_0} \arrow[rrr] \arrow[dddrrr, dotted] & & & Y_{I_0,I} \arrow[ddd, "u_{I_0,I}"] \\
 & & & Y_{I',I} \arrow[ddr, "u_{I',I}"] & \\
 & & Y_{I'} \arrow[ur] \arrow[drr, dotted] & & \\
 & & & & Y_I\ .
\end{tikzcd}
\tag*{(5~\footnotemark[9])}
\]
\footnotetext[9]{unfathomable diagram, to be replaced by diagram (6) above.}

The upper cube is the cube in $\mathcal{M}$ with 6 cocartesian faces deduced from the three arrows from $X_{I'_0}$ to $X_{I_0}$, $X_{I'}$, $Y_{I'_0}$. % [sic] faces cocartésiens -> faces cocartésiennes, gender agreement
We shall need to know that the \textit{base} of the cube is a cocartesian square. The front part of the diagram is the avatar of (2) corresponding to the inclusions $I'_0 \subset I' \subset I$, the back part the avatar for the inclusion $I'_0 \subset I_0 \subset I$. So the two other vertical squares, below the upper cube, are also cocartesian, as is the slanting dotted square, which is none other than the base $(4_Y)$ of diagram (4). From these features of diagram (5) it follows easily (although \msp{page 17} I put in a good day of hit-or-miss scribbling, before finally realizing it \ldots) that the diamond-shaped square in solid lines
\[
\begin{tikzcd}[column sep=small, row sep=small]
 & Y_{I'_0,I} \arrow[dl] \arrow[dr] & \\
Y_{I',I} \arrow[dr] & & Y_{I_0,I} \arrow[dl] \\
 & Y_I &
\end{tikzcd}
\]
is also cocartesian. We see it trivially on the unfolded diagram
\printedpage{11}
% head: A. Grothendieck : Dérivateurs, Ch. XVII, Édition des Universités de Montpellier II et Paris VII
% foot: Version 9 février 2019
\[
\begin{tikzcd}[column sep=tiny, row sep=small]
 & & Y_{I'_0} \arrow[dl] \arrow[dr] & & \\
 & Y_{I'_0,I'} \arrow[dl] \arrow[dr] & \cocart & Y_{I'_0,I_0} \arrow[dl] \arrow[dr, "u_{I'_0,I_0}"] & \\
Y_{I'} \arrow[dr] & \cocart & Y_{I'_0,I} \arrow[dl] \arrow[dr, "u^I_{I'_0,I_0}"] & \cocart & Y_{I_0} \arrow[dl] \\
 & Y_{I',I} \arrow[dr, "u_{I',I}"] & \cocart & Y_{I_0,I} \arrow[dl] & \\
 & & {?\simeq Y_I \simeq Y_{I'} \amalg_{Y_{I'_0}} Y_{I_0}\ ,} & &
\end{tikzcd} % ?? uncertain: coproduct sign -- \amalg or \sqcup
\tag*{(6)}
\]
which is finally the relevant diagram (diagram (5) was of no use, it is a wrong end of things that I took hold of). We see on diagram (6) that to prove $u_{I',I} \in C$ resp. $TC$, it \msp{page 18} suffices to prove $u_{I'_0,I_0} \in C$ resp. $TC$. But this is an assertion concerning the inclusion of ordered sets $I'_0 \subset I_0$ that are \textit{finite.} We are thus reduced to proving the following lemma, where $I_0$ becomes $I$:

\textbf{Lemma.} \textit{Let $I$ be a finite ordered set, $u_\bullet : X_\bullet \longrightarrow Y_\bullet$ an arrow in $\mathcal{M}(I^\circ)$, where $\mathcal{M}$ is a category equipped with a set $C \subset \mathrm{Fl}(\mathcal{M})$ of arrows such that $C$ is stable under cobase change and under composition, and contains the isomorphisms. We suppose $u_\bullet \in C_I$, i.e. that for every $i \in I$, the arrow $u_{I_{<i},I/i}$ is in $C$. Then for every open subset $I'$ of $I$, the arrow $u_{I',I}$ is in $C$}~\footnote[10]{\textbf{N.B.} Of course, $u_{J,K} \in C$ if $J \subset K$ are two open subsets of $I$, as we see by applying the lemma to $K$ instead of $I$.}\textit{.}

\textbf{Corollary.} \textit{If moreover $W \ \subset\ \mathrm{Fl}(\mathcal{M})$ is a mildly saturated localizer, such that $TC \overset{\textup{déf}}{=} W \cap C$ is stable under cobase change (and is so in any case under composition, and contains the isomorphisms), then we have}
\[
C_I \cap W_{I^\circ} \ = \ (TC)_I,
\]
\textit{thus if $u_\bullet \in C_I$ (cf. above) and the $u_i \in W$, then $u_\bullet \in (TC)_I$, hence applying what is said for $C$ to $TC$, we have $u_{I',I} \in TC$ for every} \msp{page 19} \textit{open subset $I'$ of $I$.}

We shall prove the lemma and its corollary at the same time, proceeding by induction on $\operatorname{card} I = n$. If $n = 0$, i.e. $I = \emptyset$, the statements are trivial. By transitivity, using diagrams of type (2), we are reduced to the case where $I = I' \cup \{i_0\}$, where $i_0 \in I \backslash I'$. We again set $I_0 = I/i_0$, $I'_0 = I_0 \cap I'$, and the reasoning we have just carried out shows us that it suffices to prove that $u_{I'_0,I_0} \in C$ resp. $\in TC$. If $I'_0 \neq I$, this follows from the induction hypothesis. % [sic] I'_0 \neq I -> I_0 \neq I, the alternative « Sinon, i_0 est un objet final de I » requires I_0 = I
Otherwise, $i_0$ is a final object of $I$, and $I' = I_{<i_0}$, and then $u_{I',I} \in C$ is included in the definition of $C_I$. Remains then to prove that $u_{I',I} \in TC$ if the $u_i \in W$. We have the diagram
\printedpage{12}
% head: A. Grothendieck : Dérivateurs, Ch. XVII, Édition des Universités de Montpellier II et Paris VII
% foot: Version 9 février 2019
\[
\begin{tikzcd}
X_{I'} \arrow[r] \arrow[d, "u_{I'}"'] \arrow[dr, phantom, "\cocart"] & X_I \arrow[d, "\tilde{u}_{I',I}"'] \arrow[dd, dotted, bend left=40, "u_I"] \\
Y_{I'} \arrow[r] & Y_{I',I} \arrow[d, "u_{I',I}"'] \\
 & Y_I\ ,
\end{tikzcd}
\]
\msp{page 20} and the induction hypothesis implies $u_{I'} \in TC$, hence $\tilde{u}_{I',I} \in TC$ by cobase change. On the other hand, the arrow $u_I : X_I \longrightarrow Y_I$ identifies with $X_{i_0} \longrightarrow Y_{i_0}$ ($i_0$ being a final element of $I$), so it is in $W$. Thus by saturation, $u_{I',I} \in W$, hence $u_{I',I} \in TC$, q.e.d.

Thus, proposition 2 is proved, as well as its corollary 1 -- and we see that if in proposition 2 and its corollary, we suppose $I$ \textit{finite,} then it is unnecessary in the hypotheses for $C$, resp. $TC = C \cap W$, to suppose stability under composition that is \textit{transfinite,} ordinary composition suffices. But it remains to prove corollary 2, which I have just proved in the case $I$ \textit{finite,} without having to appeal to a ``transfinite'' hypothesis. We must now prove that it is still OK in the locally finite case, supposing $C$ stable under transfinite composition. (There is no longer any need to return to the case of $W$, for we now already know that $C_I \cap W_{I^\circ} = (TC)_I$, by corollary 1 already proved.)

\msp{page 21}
The proof is essentially the same as for the special case already treated $I' = \emptyset$, $I'' = I$, and we could just as well have written up the proof in this setting, with the help of a strictly increasing transfinite sequence of open subsets in $I$, $I_\alpha$, with \fbox{$I_0 = I',$} $I_{\alpha_0} = I''$ $(= I)$, and
\[
\begin{aligned}
I_\alpha \ &= \ I_{\alpha-1} \cup I/s_{\alpha-1} && \text{if $\alpha$ nonlimit, $I_{\alpha-1} \neq I$, choosing $s_{\alpha-1} \in I \backslash I_{\alpha-1}$,} \\
I_\alpha \ &= \ \textstyle\bigcup_{\beta<\alpha} I_\beta && \text{if $\alpha$ limit,}
\end{aligned}
\]
and again setting
\[
X_\alpha \ = \ X_{I_\alpha}, \qquad Y_\alpha \ = \ Y_{I_\alpha}, \qquad u_\alpha = u_{I_\alpha} \ : \ X_\alpha \longrightarrow Y_\alpha.
\]
Let $\alpha_0$ be such that $I_{\alpha_0} = I$ (terminal index of the sequence), we must prove that in
\[
\begin{tikzcd}
X_0 \arrow[r] \arrow[d] & X_{\alpha_0} \arrow[d] \\
Y_0 \arrow[r] & Y_{0,\alpha_0} \arrow[d, "u_{0,\alpha_0}"] \\
 & Y_{\alpha_0}\ ,
\end{tikzcd}
\]
\printedpage{13}
% head: A. Grothendieck : Dérivateurs, Ch. XVII, Édition des Universités de Montpellier II et Paris VII
% foot: Version 9 février 2019
\msp{page 22} we have
\[
u_{0,\alpha_0} \ \in \ C.
\]

Taking into account the marginal supplement to proposition 1 (which is proved in the same way \ldots), we are reduced to proving that the homomorphism of the ordinal systems under consideration is in $C_{I_{\alpha_0+1}}$ (where $I_{\alpha_0+1}$ is the set of ordinals $\leq \alpha_0$). Now, that this is so is proved exactly as earlier, except that we must first make a verification of $(*)$ p.~1 ($\tilde{u}_\bullet \in C$) for $i_0 = 0$. Then diagram $(*)$ p.~1 is
\[
\begin{tikzcd}
\emptyset \arrow[r] \arrow[d, "\mathrm{iso}"] & X_0 = X_{I'} \arrow[d, "\mathrm{iso}"', "\mathrm{id}"] \\
\emptyset \arrow[r] & X_0 = X_{I'} \arrow[d, "u_0 = u_{I'}"] \\
 & Y_0 = Y_{I'}\ ,
\end{tikzcd}
\]
and we must therefore verify $u_0 = u_{I'} \in C$. Now this is contained in proposition 2 itself, already proved, which we must therefore use. % [sic] prouvé -> prouvée, agreement with « la proposition 2 »
The rest is proved as before. Thus, corollary 2 of proposition 2 is also proved.

\textbf{Comments:} There is an uneasiness, for proposition 1 and proposition 2 (as well as their corollaries) have essentially the same conclusion, with hypotheses on $I$ which are \msp{page 23} different: well-ordered in one case, locally finite in another. The natural idea is to give a common statement, which covers both at once, with $I$ artinian, i.e. every decreasing sequence is stationary. What is needed is a hypothesis on $I$ on the one hand, on $u_\bullet : X_\bullet \longrightarrow Y_\bullet$ in $\mathcal{M}(I^\circ)$ on the other hand, which ensures this:

\begin{enumerate}
\item[$(*)$] $\forall\ i_0 \in I$, and every open subset $J'$ of $J = I/i_0$, the arrow $u_{J',J}$ in the diagram
\[
\begin{tikzcd}
X_{J'} \arrow[r] \arrow[d, "u_{J'}"'] \arrow[dr, phantom, "\cocart"] & X_J \arrow[d, "\tilde{u}_{J',J}"'] \arrow[dd, bend left=40, "u_J"] \\
Y_{J'} \arrow[r] & Y_{J',J} \arrow[d, "u_{J',J}"'] \\
 & Y_J
\end{tikzcd}
\]
is in $C$.
\end{enumerate}

\textbf{Question:} is it correct that it suffices, for the validity of $(*)$, to suppose % [sic] exacte -> exact, agreement with « il »
\begin{enumerate}
\item[a)] $u_\bullet \in C_I$, i.e. $u_{J',J} \in C$ when $J = I/i_0$, $J' = J \backslash \{i_0\}$, and
\item[b)] \textit{$I$ locally finite?}
\end{enumerate}